\subsection{The complete first-seed lattice and its next local norm conditions}

\begin{statement}{A.165}{Proved}
{All first completed \(S\)-unit seeds, with exact finite and infinite local lifting loci}
Use the fields, embeddings and set \(S\) of A.162.  Put
\[
 \begin{gathered}
 U=\alpha^2-2,\qquad \beta_1=1-\alpha,\qquad
 g_2=\alpha-1,\quad g_{11}=2+\alpha,\quad
 g_{59}=2+2\alpha-\alpha^2,\\
 h_q=q/g_q^3\quad(q=2,11,59),\qquad
 \widehat E_i=\varprojlim_r
 \mathcal O_{K_i,S_i}^{\times}/
       (\mathcal O_{K_i,S_i}^{\times})^{3^r}.
 \end{gathered}
\tag{A.165.1}
\]
The complete fibre of \(N:\widehat E_1\to
\widehat{\mathcal O_{\Q,S}^{\times}}^{(3)}\) above \(-2\)
is parametrized bijectively by
\[
 \mathfrak b(z,a,b,c)=
       \beta_1 U^z h_2^a h_{11}^b h_{59}^c,
       \qquad(z,a,b,c)\in\Z_3^4.
\tag{A.165.2}
\]
It has a preimage in \(\widehat E_2\) precisely on the lattice
\[
                 a\in3\Z_3,\qquad b-z,c-z\in3\Z_3.
\tag{A.165.3}
\]
There is such a preimage whose localization at every place of \(K_2\)
is a completed norm from \(K_3\) precisely when
\[
        a\in9\Z_3,\qquad
        b-z-3\in9\Z_3,\qquad c-4z-6\in9\Z_3.
\tag{A.165.4}
\]
Thus the complete surviving first-seed classes modulo nine are
\[
\begin{array}{c|rrrrrrrrr}
z&0&1&2&3&4&5&6&7&8\\ \hline
a&0&0&0&0&0&0&0&0&0\\
b&3&4&5&6&7&8&0&1&2\\
c&6&1&5&0&4&8&3&7&2 .
\end{array}
\tag{A.165.5}
\]
There exists a preimage in \(\widehat E_2\) that belongs, at every
place, to a compatible completed local norm family through all higher
levels precisely on the smaller locus
\[
              a=0,\qquad
              b-z-3,\ c-4z-6\in9\Z_3.
\tag{A.165.6}
\]
These assertions specify the completed global first preimage and its
actual local norm maps.  They do not assert a completed global
\(S\)-unit preimage in \(K_3\), or an infinite global norm family.
\end{statement}

\begin{proof}
\emph{The complete group in the first field.}
The ring of integers of \(K_1\) is \(\Z[\alpha]\).
Its power-basis discriminant is \(-243\); only three can divide its
index, and at three \(\alpha\) is a uniformizer in the totally ramified
cubic extension.  Grouping the convergent uniformizer expansion into
the three powers \(1,\alpha,\alpha^2\), using \(\alpha^3=3\), shows
that its local integers are \(\Z_3[\alpha]\).
The Minkowski bound is
\[
                 \frac{8\sqrt3}{\pi}<4.42 .
\tag{A.165.7}
\]
The two primes above two are generated by
\(g_2=\alpha-1\) and \(\alpha^2+\alpha+1=2/g_2\), of norms two and four;
the three-prime is generated by \(\alpha\).
Every ideal class has a representative of norm at most four, so these
identities prove \(\operatorname{Cl}(K_1)=0\).
The remaining generators and their complements satisfy
\[
\begin{aligned}
 (2+\alpha)(4-2\alpha+\alpha^2)&=11,\\
 (2+2\alpha-\alpha^2)(10-\alpha+6\alpha^2)&=59 .
\end{aligned}
\tag{A.165.8}
\]
The first factor in each line has norm \(q\), and the second norm
\(q^2\), by A.162.4.  These are the degree-one and degree-two primes
over \(q\), since \(q\equiv2\pmod3\).

The unit rank is one.  The unit \(U\) has norm one and inverse
\(4+3\alpha+2\alpha^2\).  Under the actual residue map
\(\alpha\mapsto4\pmod{61}\), its value is fourteen, with
\[
                    14^{20}=13\ne1\pmod{61}.
\tag{A.165.9}
\]
It is therefore not a cube.  Its index in the free unit group is prime
to three: a factor of three in that index would make \(U\) a cube
up to sign, and the sign is itself a cube.
Consequently \(U\) is a basis of the pro-three unit completion.
The sign group has trivial pro-three completion.

The seven prime valuations, together with \(U\), now give the entire
completed \(S\)-unit group.  On those prime coordinates the rational
norm has rows \((1,2),(1),(1,2),(1,2)\), at \(2,3,11,59\).
Solving these four equations above the literal seed \(\beta_1\)
gives exactly A.165.2, with uniqueness.

\emph{A complete finite certificate for the second \(S\)-unit group.}
Order the primes of \(K_2\) over \(S\) as
\[
  P_0,P_2,P_6,P_3,P_{11,0},P_{11,2},P_{59,0},P_{59,2}.
\tag{A.165.10}
\]
The subscripts \(0,2\) at eleven and fifty-nine distinguish the old
degree-one and degree-two residue fields; their relative ramification
index is three.
The three literal integral elements
\(\sigma_{11},\sigma_\Sigma,\sigma_{59}\), specified in the complete
coefficient supplement for this statement, have divisors
\[
 \begin{aligned}
 \operatorname{div}(\sigma_{11})&=P_{11,2},&
 N_{K_2/\Q}\sigma_{11}&=-11^2,\\
 \operatorname{div}(\sigma_\Sigma)&=P_{11,0}+P_{59,0},&
 N_{K_2/\Q}\sigma_\Sigma&=-649,\\
 \operatorname{div}(\sigma_{59})&=P_{59,2},&
 N_{K_2/\Q}\sigma_{59}&=-59^2.
 \end{aligned}
\tag{A.165.11}
\]
Their displayed coordinates are integers in
\(e_j=t^j/649^{\lfloor2j/3\rfloor}\).  These basis elements are integral
because \(e_j^9=3^j649^{6j-9\lfloor2j/3\rfloor}\) is an integer.
The norm and local residue calculations in the retained certificate
prove A.165.11.  In particular the prime norm factorization excludes
every other prime from their divisors.

Set
\[
\begin{gathered}
 s_1=\beta_2,\quad s_2=\beta_1/\beta_2,\quad
 s_3=2/(\alpha-1),\quad
 s_4=\frac{t}{(\sigma_{11}\sigma_\Sigma\sigma_{59})^2},\\
 s_5=2+\alpha,\qquad
 s_6=\sigma_{11},\quad s_7=\sigma_\Sigma,\quad s_8=\sigma_{59}.
\end{gathered}
\tag{A.165.12}
\]
Their valuation vectors in A.165.10 are
\[
 e_1,\ e_2,\ e_3,\ e_4,\ 3e_5,\ e_6,\ e_5+e_7,\ e_8.
\tag{A.165.13}
\]
For \(s_4\), use the literal equation \(t^3=649^2\alpha\), which gives
\(\operatorname{div}(t)=P_3+
2(P_{11,0}+P_{11,2}+P_{59,0}+P_{59,2})\).

Together with the four units \(u_i\) of A.159, these eight elements
give a basis of \(\widehat E_2\).
Here is the finite independence verification behind that assertion.
Evaluate the four units and the eight literal elements A.165.12 at
the following degree-one primes, in that order:
\[
\begin{array}{c|r|rrrrrrrrrrrr}
p&t&\multicolumn{12}{c}{\text{cubic residue character rows}}\\ \hline
2029&145&2&2&0&2&2&1&1&2&1&0&1&2\\
2029&510&2&2&0&2&0&0&1&2&1&2&0&2\\
2029&1374&2&2&0&2&1&2&1&2&1&1&0&1\\
2131&1155&2&2&1&0&0&2&2&1&2&1&0&2\\
2131&1397&2&2&1&0&1&1&2&0&2&2&0&1\\
2131&1710&2&2&1&0&1&1&2&1&2&1&2&2\\
2203&1250&0&2&1&1&2&1&0&0&2&0&1&1\\
2203&1589&0&2&1&1&1&2&0&0&2&0&1&0\\
2221&529&0&1&0&1&2&0&2&0&2&0&0&0\\
2251&108&0&2&0&1&0&0&0&1&1&1&2&1\\
2251&258&0&1&1&2&0&2&1&2&2&1&0&2\\
2383&911&2&0&0&1&0&0&0&1&0&0&0&2
\end{array}
\tag{A.165.14}
\]
The matrix has rank twelve over \(\F_3\).  Every evaluation is at the
literal root \(t\) shown in its row; the chosen nonidentity cubic root
and all intermediate values are retained in the finite certificate.
All twelve elements are actual \(S\)-units by A.159 and
A.165.11--13.
Dirichlet's theorem gives \(S\)-unit rank \(4+8=12\).
The rank matrix first proves multiplicative independence, and then
excludes a three-factor in the subgroup index by exactly the cube-root
argument of A.159.5.  The four units already span the completed global
unit group.  This proves the completed basis assertion directly,
without any class-group or regulator completeness assumption.

\emph{The exact primitive norm direction.}
Write
\[
 w_0=\frac{\sigma_{11}\sigma_{59}}{\sigma_\Sigma^2}.
\]
Its relative norm has the same divisor as \(h_{11}h_{59}\).
Hence there is a unique \(k\in\Z_3\) with
\[
                    N_{K_2/K_1}w_0=U^k h_{11}h_{59}
                    \quad\text{in }\widehat E_1.
\tag{A.165.15}
\]
The sign disappears in this pro-three identity.
The exact finite calculations for the shorter representatives give
\[
 \begin{gathered}
 k\equiv7\pmod9,\\
 (\theta_{11}(w_0),\theta_{59}(w_0))=(1,0),\qquad
 (\theta_{11}(U),\theta_{59}(U))=(2,2).
 \end{gathered}
\tag{A.165.16}
\]
Here \(\theta_q\) is the additive stable norm character, with generator
\(\lambda_q(\alpha)\), for the actual extension based on \(\rho_2\).
More explicitly, if \(x=\pi^v c\) and \(\pi=t^2/q\), then
\[
 \theta_q(x)=
 \log_{\lambda_q(\alpha)}\lambda_q(c)+v\pmod3.
\tag{A.165.17}
\]
Indeed \(\rho_2/\pi^8=(649/q)^{-8}\alpha^{-5}\) has character
\(\lambda_q(\alpha)\), and \(8^{-1}\equiv2\pmod3\) in A.162.9.
The order-nine residue checks proving \(k\) use a primitive character
of \(U\) at degree-one primes of \(K_1\), after dividing by the exact
prime generators in A.165.8.  All powers and relative norm
polynomials are retained; the complete coefficient products are not
expanded during the residue calculation.

Since \(k\equiv1\pmod3\), the following element is defined in the
completed \(S\)-unit group:
\[
 \eta=w_0 U^{(1-k)/3},\qquad
 N\eta=U h_{11}h_{59},\qquad
 (\theta_{11}(\eta),\theta_{59}(\eta))=(0,2).
\tag{A.165.18}
\]
The last identity follows from
\((1-k)/3\equiv1\pmod3\) and A.165.16.
Both \(w_0\) and \(U\) have zero valuation at every place above two.
No assertion that \(\eta\) is an actual field element is needed:
continuous powers by \(\Z_3\) are operations in the stated completion.

\emph{The complete first norm image.}
Necessity of A.165.3 can be read from the actual local maps.
At the old degree-two prime over two, \(K_2/K_1\) is unramified cubic,
so its norm valuation is divisible by three; this is \(a\in3\Z_3\).
At the quadratic places over eleven and fifty-nine,
\(\rho_1=649^8\alpha\) has valuation eight and unit character
\(\lambda_q(\alpha)\).  The residue computations give
\[
 \lambda_q(\beta_1)=1,\qquad
 \lambda_q(U)=\lambda_q(\alpha)^2.
\]
Since \(8^{-1}\equiv2\pmod3\), the literal stable-character
equations are \(2z-2b=0\) at eleven and \(2z-2c=0\) at
fifty-nine, modulo three.
These are the remaining two congruences in A.165.3.
They apply to completed norms because the local residue quotient is
finite and the completed local norm image is compact.

Conversely, on A.165.3 the literal completed preimage
\[
 B_2(z,a,b,c)=
 \beta_2\,\eta^z\,
 h_2^{a/3}h_{11}^{(b-z)/3}h_{59}^{(c-z)/3}
\tag{A.165.19}
\]
has norm A.165.2, by A.165.18.
This proves the exact global \(S\)-unit norm image.

\emph{The next local norm test and its invariance across the fibre.}
The two stable characters of \(\beta_2\) are \((2,1)\);
those of \(h_2\) are zero; \(h_{11}\) contributes \((1,0)\)
and \(h_{59}\) contributes \((0,1)\).  All these are the literal
uniformizer comparisons of A.162.7--11.
Thus A.165.19 has characters
\[
 \theta_{11}(B_2)=2+\frac{b-z}{3},\qquad
 \theta_{59}(B_2)=1+2z+\frac{c-z}{3}
                         \quad\text{in }\F_3.
\tag{A.165.20}
\]
Their simultaneous vanishing is exactly the last two congruences
of A.165.4.
Every other completed \(S\)-unit preimage differs by the relative
norm-one group of A.159.12.  That group has valuation zero at both
changed primes and trivial unit characters there.  Hence A.165.20
is independent of the chosen \(S\)-unit preimage.

At two, A.165.19 has divisor
\[
       (1-2a/3)P_0-(2a/3)P_2+(a/3)P_6.
\tag{A.165.21}
\]
The \(P_2\) and \(P_6\) branches become unramified cubics in
\(K_3/K_2\).  A.159.10--12 show that every relative norm-one
\(S\)-unit has divisor a multiple of \(P_2-2P_0\); it cannot change
the \(P_6\) valuation.  Consequently \(a/3\in3\Z_3\) is necessary
for every preimage.  Under this condition the displayed preimage
already has both unramified valuations divisible by three.
At the distinguished branch the degree-one factor supplies all norms.
At three A.140 gives the onto completed norm map.  At the linear
eleven/fifty-nine branches every residue unit is a cube.  Outside
these places the target is a unit in an unramified root algebra.
The complex factors and the chosen real factor give the archimedean
norm maps directly.  This proves both directions of A.165.4.

\emph{All heights on the same fixed local towers.}
At \(P_6\) the residue of \(\rho_2\) is a primitive ninth root of
unity.  Its \(3^{m-2}\)-th roots generate the unramified extension
of relative degree \(3^{m-2}\), because
\(\operatorname{ord}_{3^m}(2)=2\cdot3^{m-1}\).
At \(P_2\) the analogous relative residue degree is the same, starting
with a primitive third root.  The invariant \(P_6\) valuation
\(a/3\) must therefore lie in \(3^{m-2}\Z_3\) for every \(m\).
This is precisely \(a=0\).
The two tame characters in A.165.20 remain the full norm restrictions
at every positive depth, since \(v_3(11^2-1)=v_3(59^2-1)=1\).
These observations prove necessity of A.165.6.

Conversely on A.165.6, the same element A.165.19 has zero valuation
at \(P_2,P_6\), and its valuation at \(P_0\) is one.
The fixed degree-one factor handles \(P_0\); unramified unit norms
handle the other two branches.  Its two stable characters vanish
at every higher tame depth, and the three-adic norm maps are onto.
The remaining unit and archimedean arguments are the same as above.
Completed local multiplicative groups are compact at every finite
place.  At each finite top height a norm preimage therefore supplies
all lower compatibility equations; the closed fibres have the finite
intersection property.  Their product gives an actual compatible
completed local family.  At the real and complex places the choices
of factors already give compatible maps explicitly.
This proves A.165.6, with exactly its stated local scope.
\end{proof}

\paragraph{Reproducible finite data.}
The first-field coordinates and both affine congruence reductions are
checked in \path{code/first_seed_sunit_fibre_certificate.py}.
The final literal representatives and their rank-twelve matrix are
checked in
\path{code/two_inert_short_sunit_principal_certificate.py} and
\path{code/two_inert_short_sunit_saturation_certificate.py};
the complete printed coefficient blocks are independently reconstructed
in \path{code/two_inert_short_sunit_table_independent_audit.py}.
An independent discovery-product verification is retained through the
complete compact support, S-valuation, tame-angular-component,
rank-twelve and order-nine norm certificates
\path{code/two_inert_compact_support_certificate.py},
\path{code/two_inert_compact_sunit_s_valuation_certificate.py},
\path{code/two_inert_compact_sunit_tame_certificate.py},
\path{code/two_inert_compact_sunit_saturation_search.py}, and
\path{code/two_inert_first_norm_image_certificate.py}.
All inputs, full integer coordinates and executed outputs are retained.
The compact products and the shorter representatives have identical
divisors; their two separate normalizations give the same exact
\(\eta\)-characters in A.165.18.

\subsection*{Literal principal generators used in A.165}
Put $B=10^{40}$ and $e_i=t^i/d^{\lfloor2i/3\rfloor}$ for $0\leq i\leq8$.
The three elements below are defined by the displayed integer coordinates:
\[\sigma_j=\sum_{i=0}^8 a_{j,i}e_i,\qquad
a_{j,i}=\epsilon_{j,i}\sum_{k\geq0}b_{j,i,k}B^k.
\]
Each row supplies two successive forty-digit blocks; unlisted blocks are zero.
The sign is repeated on each row.  Every $e_i$ is integral because
$e_i^9=3^i d^{6i-9\lfloor2i/3\rfloor}\in\Z$.
Their exact rational norms are $-121,-649,-3481$, respectively,
and their exact prime divisors are certified in A.165.
\begingroup\setlength{\tabcolsep}{3pt}\fontsize{7.6}{9.2}\selectfont\ttfamily
\par\medskip\normalfont\normalsize Generator $\sigma_{11}$\par\smallskip
\fontsize{7.6}{9.2}\selectfont\ttfamily
\begin{longtable}{rrrll}
$i$ & $k$ & $\epsilon$ & $b_{j,i,k}$ & $b_{j,i,k+1}$ \\ \hline\endhead
0 & 0 & + & 4085765469006991614846224650275985123879 & 2807453240084609810433909756602982388463 \\
0 & 2 & + & 1308421798136464620933797602795618501789 & 6750264511899386519826072800169773257540 \\
0 & 4 & + & 5827484447407566631582697086872722670576 & 2593328686421282198929886108464317659777 \\
0 & 6 & + & 9753280323425001967959757951095029033147 & 1140921001632188353816612984250638310384 \\
0 & 8 & + & 6904340104205139748440817532734067368038 & 7607177367698926752327232627274357739611 \\
0 & 10 & + & 2151101335322826748999713218517578059325 & 9733133071683241297290868540123936458167 \\
0 & 12 & + & 2390653906548344049403246568800655756927 & 8751955228265581316125508539561022871049 \\
0 & 14 & + & 4706679000005614080627892601897622450621 & 7632742321792417912651358008861791235889 \\
0 & 16 & + & 2575243044620015909791144382225374246955 & 2497642224374485201251662862936912473388 \\
0 & 18 & + & 6497403135825525945992081836505805941420 & 4750085739865957143423001625786297536757 \\
0 & 20 & + & 4751640386720365178759718236621536323983 & 1522161506915964779941442380924272845256 \\
0 & 22 & + & 5528691483236264288865060903561761873710 & 0407574768839544458657451541899014976367 \\
0 & 24 & + & 8945907733294315954084712275157594034025 & 2230045873313416440207587962158512881977 \\
0 & 26 & + & 2179012604013386242497204453390498740852 & 2646624271931557896845480820748998966938 \\
0 & 28 & + & 2927762390264260481261864088044077683500 & 0466202321423110221324031881036299115644 \\
0 & 30 & + & 3772026080315622772984446323194549143825 & 9846467429132831444754118314696152354934 \\
0 & 32 & + & 2025120058685606025385684475231077137419 & 7939651289607622050675560067921482276439 \\
0 & 34 & + & 3017775328581744723673941197968967193708 & 0595860175350580196074323239905351329233 \\
0 & 36 & + & 2406158462699546147074543287906786439194 & 0000000000000000000000000000000535588091 \\
1 & 0 & + & 6022137300505656763518396236115489735897 & 3070382753935186539047023970562576995172 \\
1 & 2 & + & 7209078244844311808025757073778324853936 & 6630565227294913940648426935160669947479 \\
1 & 4 & + & 5228415802772012441636492653903416993026 & 9026313276894350303071275290594290217307 \\
1 & 6 & + & 5904731563772406840423045201348404687469 & 6947153771113209088081285094629931718126 \\
1 & 8 & + & 2989067358372605300533226244641641778196 & 5092646229370287395646151139237639264504 \\
1 & 10 & + & 6722092981595518657977423904812738106282 & 1056528058580969651259386325007199773564 \\
1 & 12 & + & 9867349686618174734924392670731389506218 & 7665176655614217192624283990330548320989 \\
1 & 14 & + & 8143813935078198794586353924313345675244 & 8147241723532505547159966433983120665786 \\
1 & 16 & + & 3639942633419950645539254909833344488676 & 5247411271711826304344080439315095577805 \\
1 & 18 & + & 8603449587920499750063100299959460228211 & 5894920721945369185532319975647548883705 \\
1 & 20 & + & 0965817123582100839148243607949970360772 & 7520910814433925337258575113524797914810 \\
1 & 22 & + & 3344735428745427224166964165693387224872 & 4017339354053829962220818023226629806536 \\
1 & 24 & + & 1188990617468937110926440437232926633042 & 1291510773862130581622517545044202273089 \\
1 & 26 & + & 2888015910503357605779912669777021740263 & 5073871169783528153135768831378496907000 \\
1 & 28 & + & 4746571719401944671766395570631189065686 & 5777197327147714825849517169015272346337 \\
1 & 30 & + & 4706079926011731081390322189476208296803 & 9315017027860960011981031079097239390333 \\
1 & 32 & + & 1784285756802339420098570730562155800064 & 5543208743432378256931999772882073308204 \\
1 & 34 & + & 4060624573103207378014169431700855143725 & 3400571204615032192515721049361643172184 \\
1 & 36 & + & 1109070286576713024078066008892566111433 & 0000000000000000000000000000000009779719 \\
2 & 0 & - & 0130962421063106224328215887626724738757 & 9277628222219048022689450885025637806759 \\
2 & 2 & - & 1580468742089055087646043331430415007264 & 8285386025847961943852346050834865998952 \\
2 & 4 & - & 0378102880752717615519546983856752049202 & 5636099776250016113883430192083024983660 \\
2 & 6 & - & 6274402008509325145826164688767567586228 & 7975927331701071252216257341838371771565 \\
2 & 8 & - & 3561615231557703839094240114661270831883 & 4010547148478408967949845244594032086894 \\
2 & 10 & - & 0193100652113097134477526538990270712324 & 1929801814374117353584955090286623448797 \\
2 & 12 & - & 5170366103720185315351436335889249693224 & 3198781658391882131596119285839915168204 \\
2 & 14 & - & 3323396532967771298286883483426429318212 & 2394432173346603580775695207701632160165 \\
2 & 16 & - & 6935180207149629825156604783310312208942 & 0481535403148259957771438731486257932101 \\
2 & 18 & - & 3347409062188358290870885174460131255354 & 5469206898905647937808781098702851977324 \\
2 & 20 & - & 0866814854454277500019649486797392854891 & 1396466343508596000385119626946813866470 \\
2 & 22 & - & 3408632943602181716602956980105139172497 & 3953112545025185473483627000093347801352 \\
2 & 24 & - & 0563228280117627247814422667513896529540 & 2158588388880186923914796587247546856456 \\
2 & 26 & - & 2449515284626760379850674537842488369884 & 3501363915931188600786973523120193154035 \\
2 & 28 & - & 6058697315196445317389559485360025133837 & 7604589879199578787216926319058644960112 \\
2 & 30 & - & 7196194786650609348026175405921113178309 & 5555279436502019769651370349070795289210 \\
2 & 32 & - & 4056940178827518621797849474876909445061 & 3550083739348479444941010544774350304254 \\
2 & 34 & - & 2518535793964469120536030337328408112731 & 9360094167656827665501554651907737560708 \\
2 & 36 & - & 3188935180328701661831744567372873403075 & 0000000000000000000000000000000022433396 \\
3 & 0 & - & 8289709628814897542797395329814040774624 & 5734376825845060176896335317256383261763 \\
3 & 2 & - & 8720803715465061283079229835535378793437 & 8380159707992683326878982086428310937421 \\
3 & 4 & - & 7783888649238525914298032937573781348309 & 1128605357825136524769017585809017674359 \\
3 & 6 & - & 9080955749838134153239193781448067057055 & 3737654871706970972867299308788810266796 \\
3 & 8 & - & 5633810536705927767445542234543010259323 & 8097506767371226868568930878007857392540 \\
3 & 10 & - & 0994145191428849207135167082104537919524 & 9761376915099894121359154681378293565215 \\
3 & 12 & - & 1906749596400555372746481472608027034917 & 9651378616816736534267603373414622117967 \\
3 & 14 & - & 6041570853993547819391110103278885728164 & 4461807872170235919321810751028624498015 \\
3 & 16 & - & 0188987682081764434369732503952229870031 & 6345892441789178113678958445134457043455 \\
3 & 18 & - & 6161375934318000773255918399901870950671 & 4770682419098126244078034313059932334613 \\
3 & 20 & - & 6556227541328851217070423265993652426451 & 7852275941671326826730036898275814992857 \\
3 & 22 & - & 8969227153984931600332187494670757633293 & 2996416347732821298709228447675391208773 \\
3 & 24 & - & 1617959411820060599019240967181595773508 & 1256406155920234929852099940340464251055 \\
3 & 26 & - & 2045523464021793825269974437401249351096 & 5506162703800722420525380843756107244437 \\
3 & 28 & - & 1205676809676113145640339504666645145947 & 5392143397001017610571158026767004986080 \\
3 & 30 & - & 0977480020986980121129898410402962557807 & 1441242557936463294609842311807547809515 \\
3 & 32 & - & 4507204082000216824444882929967818529608 & 6064553870536415741347706835541698738246 \\
3 & 34 & - & 3926963310569277791709454588010566834764 & 2008943845217199438765885953416219284590 \\
3 & 36 & - & 7444571637531516187697362452728569495967 & 0000000000000000000000000000000633990400 \\
4 & 0 & - & 9191835734668153326965419330265888982324 & 7682778368108080609341235526754443643248 \\
4 & 2 & - & 8229705096335131447926009735786955134602 & 9683374547173186963921868893899650607466 \\
4 & 4 & - & 4856459115808921860217752977459959529085 & 2657321885086397991564655301270420861499 \\
4 & 6 & - & 6382450855484208792458714511985705138211 & 4176791084558671132358585985644945073011 \\
4 & 8 & - & 7655836642969691748692645893809284093938 & 3670235190911292675908979262868030919585 \\
4 & 10 & - & 7318151257678710342329158277635268923487 & 8317426869198969695962372340947709587225 \\
4 & 12 & - & 9380957957239771569777387821033993233185 & 5283418102731452316456058697784214712497 \\
4 & 14 & - & 4455454274416969573184086861215424197147 & 5202115028669549283245868100561936101419 \\
4 & 16 & - & 0792812021223051094203668732401196929643 & 8818223936095931663086741337097193366490 \\
4 & 18 & - & 3107207433998177376582999813540417892933 & 5829861357558761378503396168129336217359 \\
4 & 20 & - & 9850737325523889674323549146421092511534 & 9460362910890269189100218338803368366540 \\
4 & 22 & - & 8358914559068801286819525267495537018151 & 5645394187138325434962014459005957240025 \\
4 & 24 & - & 2430016280534476051178047846877452051534 & 3706751862467484791093426738306104500080 \\
4 & 26 & - & 4963668047144201827257645296794256338893 & 3919861401754219567613500472255374229502 \\
4 & 28 & - & 2958013576603471946584772281761711673727 & 8243054312649028408283407671529764261356 \\
4 & 30 & - & 6808568651266210071994823429774148160875 & 2328705005614606258643303877808228866994 \\
4 & 32 & - & 4361123267489956342974090102806841519274 & 9005138975553472754967939421436876026000 \\
4 & 34 & - & 3709623656527320808499333194397602077729 & 9489302052709055962949319912647188455043 \\
4 & 36 & - & 4895525248172120426598636709541161869181 & 0000000000000000000000000000000000569998 \\
5 & 0 & + & 4731079794203140096826286824515521544780 & 9864413951093871204259814845952691186744 \\
5 & 2 & + & 5642388764489631424901840373178572741446 & 8363512764169195377458832405543206577192 \\
5 & 4 & + & 0495305151582386064302278317979121058187 & 7994728587165193910469518696825386077523 \\
5 & 6 & + & 6502604339194978834894238189787586787396 & 8012518571571698296947604120428160786055 \\
5 & 8 & + & 0252828506787155988145681626939952669177 & 3078458333211143280145321402971325165428 \\
5 & 10 & + & 7286531360225235825835454254609703868375 & 2611182789047238139516054768144851070300 \\
5 & 12 & + & 7830587545357259403120974311215740244459 & 1713517250992273891638626489589948984786 \\
5 & 14 & + & 8875330776861564996757856946391385456580 & 3107695915574170832887452329243980051736 \\
5 & 16 & + & 2561066695373734656567634977429200964056 & 0038845884744908614042025761236387759614 \\
5 & 18 & + & 4556320412307541819891793661137674255687 & 1606237044524083491813830835253469235895 \\
5 & 20 & + & 3637826681181330889800309905676417483977 & 5996565745529650157424927703050074408961 \\
5 & 22 & + & 4467253728135367019075347218996424532433 & 3326283138772592842380793876881319319845 \\
5 & 24 & + & 4301395258344445905006495079729574078998 & 1179098430429447670967371728862701826143 \\
5 & 26 & + & 3182368069375947528465963307581105292566 & 6132059336720344597622097908716684198451 \\
5 & 28 & + & 1501274347300806434145235730650530171827 & 7215023278271750630804927129256490535587 \\
5 & 30 & + & 9093409959075035590539675033118172383636 & 5470941054634412715547352000460501719909 \\
5 & 32 & + & 6242836192270441638074365616972856108200 & 1891440832549827387193727431046425517946 \\
5 & 34 & + & 9310307993081503022767895126393428957229 & 4208873241440941820417620407448011411741 \\
5 & 36 & + & 0478191051667589880670937656047203038052 & 0000000000000000000000000000000055847222 \\
6 & 0 & + & 0619096666534264894329606276220033204823 & 5800713715914907236888840412471486476317 \\
6 & 2 & + & 8170221918808382919218752091311773626219 & 7985309389189988710828301325682844637229 \\
6 & 4 & + & 0278716614429004590824309912514304570828 & 0083798013618012449310843389786988397897 \\
6 & 6 & + & 3601971915064564814330591426672805718359 & 2565820082184843353612541903494372118667 \\
6 & 8 & + & 0633474130722271199845659336519113110891 & 4809020665124719067443739314058327325780 \\
6 & 10 & + & 0662187589440494021852477957846692103966 & 8664740182361172325214462155718445929276 \\
6 & 12 & + & 6348155454015143937969318940254403721197 & 5993708325319483298896304416500823670048 \\
6 & 14 & + & 0035566351806719260081218648518534663477 & 1868818808922783044163816918217346248199 \\
6 & 16 & + & 1054637080628177477603769077910266190944 & 2110562314066587630526536436846295064892 \\
6 & 18 & + & 6713438913165455482265801969291603591759 & 4262792523383648045349834814349649723610 \\
6 & 20 & + & 8134487649395144143039131464499305820870 & 8250996788801070570636319816044497979450 \\
6 & 22 & + & 4448918310829263967550618982107892013785 & 6068610717632281995064944998927547982605 \\
6 & 24 & + & 5162383534129627419715694480609096738450 & 6346624722950104413743395817450508166236 \\
6 & 26 & + & 9200833868485928102810567314910295944988 & 4461209028353672758448851727395867191847 \\
6 & 28 & + & 2352098706884995980514197608807364986215 & 6766638074154037282308669775005799813910 \\
6 & 30 & + & 5040089257741314387279663552245399263468 & 7605765325393410971477267270739492335623 \\
6 & 32 & + & 7878127545780185864822491262165630919471 & 4549156483774689077337749150546315856539 \\
6 & 34 & + & 6240470610104273938491005813394369908056 & 6264552505765092629411342216064919684824 \\
6 & 36 & + & 0350653246911832298498059400680779534114 & 0000000000000000000000000000000182100494 \\
7 & 0 & - & 5564779881105886187437418531282426377162 & 2360468310323643312096986215706267506918 \\
7 & 2 & - & 2574202940461890149009776858291101838416 & 5630402122907117650484297341189488700750 \\
7 & 4 & - & 4397714377083281357991321437547721015233 & 6060004637911836637916601897326959539640 \\
7 & 6 & - & 9574651616184555074094199913864402794775 & 6764449170524568179695814727048819687435 \\
7 & 8 & - & 3415208217625577182113653119116799930406 & 3007175709840054487795740028289587624344 \\
7 & 10 & - & 4916924905944236887849533671752539284383 & 9630807041362078351607285863500828691058 \\
7 & 12 & - & 7697151755616920467235437636240187778659 & 9276286994307124671463970914648679810526 \\
7 & 14 & - & 6932889579015482616021844110467392727926 & 0879344733441154593064277956672877968182 \\
7 & 16 & - & 8856549967452752585143770476409310744162 & 9113731362203736875539930852278082709031 \\
7 & 18 & - & 5751720036164964924779415836802054164938 & 7408357577803719155411533982619201202978 \\
7 & 20 & - & 6176837509679028231498666637264816229219 & 5325574547751662332203054224059715578069 \\
7 & 22 & - & 9519711511010959683081316660157240643953 & 4783527993935538837440603409019319161867 \\
7 & 24 & - & 2090371251278819238843813681069183811610 & 6996328395861127269406614496350794391681 \\
7 & 26 & - & 5328493592691321648587289527754764104823 & 8472640341668218227182481710436551071563 \\
7 & 28 & - & 0306354569856855653548246494420787560884 & 9820986945265358168302273582451417212275 \\
7 & 30 & - & 0643481960109736665191817535070846524730 & 5880604413839689387339844564869562759661 \\
7 & 32 & - & 9548371077179991331348486932217674690780 & 6337692409499140034534984738565015586844 \\
7 & 34 & - & 3646219535093994093787682657853449752920 & 6234300819062003177400867638382291483333 \\
7 & 36 & - & 6827168617916585193739625915582002309021 & 0000000000000000000000000000000004306383 \\
8 & 0 & - & 8271368175558610200167448923358516142667 & 7258711399019416261743989954384923841058 \\
8 & 2 & - & 4898602551923567680922685404066493576647 & 2283892837617392849759492070353850631184 \\
8 & 4 & - & 7857746392163992946656810700243764774849 & 2773599297266649812958685840550102221132 \\
8 & 6 & - & 5977490876593547180009650564295514176698 & 6672441130077472721431204609528492370929 \\
8 & 8 & - & 7753969725491427540259910688739645421472 & 8145752292797480825340157579862235342619 \\
8 & 10 & - & 4586290969721706047194716235126169381770 & 2824827112127499478048324848584796741182 \\
8 & 12 & - & 1576722224483867137490877409239533875007 & 0742970634456734402594990257548696984433 \\
8 & 14 & - & 3724393136771187208345752958048426700128 & 3764502246347145605606036541824929861214 \\
8 & 16 & - & 4138119658836616013908270160587532069637 & 9298589503899537833675731617916107177956 \\
8 & 18 & - & 2193590059946421188693193496648577763987 & 2714190948509346670974184524772821144141 \\
8 & 20 & - & 1374193025184062271186160365586091999491 & 3696593575781198734826891520011461102475 \\
8 & 22 & - & 7680359594556898686493787778366708841686 & 8694144311934021408963205814766626822824 \\
8 & 24 & - & 7675835552727592966124076080215897976408 & 2186657278336524375565363868678865900035 \\
8 & 26 & - & 8714455288696802945447847988018275154983 & 6935017965951452959998121816511418482544 \\
8 & 28 & - & 5312173954546761532778812022139138969028 & 6370341188408189761427059679289042384786 \\
8 & 30 & - & 4102469528523202111795419190744741929500 & 0467548889618481454626063738996607709899 \\
8 & 32 & - & 6472522107593618001603482627545671444427 & 4318664285855782396306428751034987377006 \\
8 & 34 & - & 4887466673394469783985771896466619599135 & 3043647618222691349097238980860821832146 \\
8 & 36 & - & 9808653934553337035742293072633317214543 & 0000000000000000000000000000000027937448 \\
\end{longtable}
\par\medskip\normalfont\normalsize Generator $\sigma_{\Sigma}$\par\smallskip
\fontsize{7.6}{9.2}\selectfont\ttfamily
\begin{longtable}{rrrll}
$i$ & $k$ & $\epsilon$ & $b_{j,i,k}$ & $b_{j,i,k+1}$ \\ \hline\endhead
0 & 0 & - & 3657358675969008555999921721308815050044 & 9432150062901437883444779929338312106675 \\
0 & 2 & - & 2761058723206255554540148437670118701179 & 0570035293181706082271627458631115189901 \\
0 & 4 & - & 1297474868585394285715691445336393275206 & 4652685438911984161749267421432920686902 \\
0 & 6 & - & 0015246138189927589589701398481842067548 & 5424937077493527597766113547382567388208 \\
0 & 8 & - & 5835931230703798658884981859375928467640 & 3791929359494636453535852240200595011021 \\
0 & 10 & - & 2496917149552811897481398089146995220211 & 7257505363014521936261572874979233404701 \\
0 & 12 & - & 4524936986741802532765886357801090142989 & 3537829015817981917503703961986103020037 \\
0 & 14 & - & 7242284576412737667543197693127979612439 & 5475884902055622010498534106800219585648 \\
0 & 16 & - & 9102239427862233124414577614157518908320 & 8219941393692203847089465425669020995306 \\
0 & 18 & - & 2597938090678799373869526757068612570241 & 7954847199949810798850089168861858492110 \\
0 & 20 & - & 8450027921723123214320487469891354427281 & 7796786532784994152640915420718749227940 \\
0 & 22 & - & 5100746439596126803184655764128836821202 & 7203555668197251606933981948780797437074 \\
0 & 24 & - & 1743454898978921930939771434098152076036 & 9042868040378606525684973480026567588798 \\
0 & 26 & - & 4863126054863754470228381040167014336062 & 7935551194548022724066864009643926132370 \\
0 & 28 & - & 2499518027913039806127188930544192630696 & 2922840171099128288926590941209116141767 \\
0 & 30 & - & 7188018931051836642278057844894949233122 & 8537762826668264644025806876672375314017 \\
0 & 32 & - & 5134975121544877344790397907716062730672 & 7612104827638206079155733317216312351666 \\
0 & 34 & - & 5594653290767981186185202700047130831986 & 9205111676265390965352318053578091219641 \\
0 & 36 & - & 6384251328200178224214882024484406478053 & 1003132967878686911649304951480402352944 \\
0 & 38 & - & 6165592629468307503197321132912231700691 & 1997458685395153323399256987984995567515 \\
0 & 40 & - & 5495577198804033143067623365511419745026 & 0000000000000000000000000000000004367850 \\
1 & 0 & - & 4130035546119668837706898533589274094650 & 0751029766749140157068748186267444983813 \\
1 & 2 & - & 0026584285371712137145669116762994716988 & 5222508209924513674435397084887392259331 \\
1 & 4 & - & 6772024144624489916390527270844668277560 & 0203407173163535202841064320528188588602 \\
1 & 6 & - & 4345135044402343546909876180959546964658 & 3362610754085266658166675626369560191976 \\
1 & 8 & - & 8422401050332402178399787353880379411068 & 6178656931820868639037725525014063394039 \\
1 & 10 & - & 6495480307340258957388787023505160725756 & 2658657221786287241641401275528486287181 \\
1 & 12 & - & 9405123802828619140192994656028101362982 & 5208090022734317563861550671893996800731 \\
1 & 14 & - & 8723486905308905011507014799325371422759 & 3259368263902195791147384691752608814462 \\
1 & 16 & - & 3876353023953202109482289251388716064237 & 0181476392546064498101642626978417687552 \\
1 & 18 & - & 1878595912870048517198989411370961833355 & 2104822330348796196960533004348032732076 \\
1 & 20 & - & 7608567288914644028541754853551388782536 & 1451788051105777723052953146194700715197 \\
1 & 22 & - & 3206204149284358443349925307274614032030 & 3618294327914081345969033909420631114175 \\
1 & 24 & - & 6564590895935513962383594703503866448273 & 5888065935614096436391516579749653788552 \\
1 & 26 & - & 9933174903029117681259198076579463197241 & 5139197629783645668238058059220270218379 \\
1 & 28 & - & 3283755999182462284360597123416898686591 & 6355755810992022476950069685329435273783 \\
1 & 30 & - & 9142533841157876849693566336180314938134 & 0692799016031346707653460845822738100933 \\
1 & 32 & - & 8600011224086183821762664044218519340996 & 0237141280333362378327703461885443604773 \\
1 & 34 & - & 5283346463131774711385183221262887586808 & 6697053146644566517527421731315265048012 \\
1 & 36 & - & 2905932709008750433849952184623189987504 & 0881213009210642288979324968810453243416 \\
1 & 38 & - & 0761985215471549088657706613000322502102 & 6967663410645488736706731864473466395708 \\
1 & 40 & - & 2499091015777400369907403916125669811912 & 0000000000000000000000000000000000051573 \\
2 & 0 & - & 3069894409233028202687982240123695515432 & 9979273201602730886047505385972566412540 \\
2 & 2 & - & 2141594320463746217510065199311244061466 & 3411421349612248885321555454746911843570 \\
2 & 4 & - & 0473509331806205191083458512074788964858 & 9261435201961591121640682493880291325639 \\
2 & 6 & - & 4445049121193392823176807755192185256298 & 4584048976150846254381484374266987340942 \\
2 & 8 & - & 9860974937612978018340804813503097780245 & 2239859146990318427827841510310205098335 \\
2 & 10 & - & 9297452875166286794896342676979403866831 & 8305097653466859166028375109669141477720 \\
2 & 12 & - & 7467096346014483428510863778546575677575 & 2743696204231764944856349482871279345688 \\
2 & 14 & - & 3075745204504109140526864225248107138139 & 9606855097918557668382266931486816791028 \\
2 & 16 & - & 9530151198012397081014255099716354973559 & 4501009689016092858629315969853631531235 \\
2 & 18 & - & 2878959114629717289766583408601930667889 & 0199941465371267296045681775591040851634 \\
2 & 20 & - & 5918322942541489658273626571262551920148 & 4043793690147307344488949731224832495403 \\
2 & 22 & - & 8857891922980617659808281299397683963713 & 8046748968860053812559338500796284157618 \\
2 & 24 & - & 8336409444021771242559450752603390674933 & 8659938231367586781664014407316258315441 \\
2 & 26 & - & 8587557189310459893613198424674531370724 & 5775476718593808077272027879318186599947 \\
2 & 28 & - & 2346490967084681818039055401066849894587 & 0324952054018307063787853567729789813859 \\
2 & 30 & - & 1591039927298030275771880373418751027338 & 1070683369222458142579910879569237314534 \\
2 & 32 & - & 8043651274528674454856686145544506804004 & 8955558043623518084918325564164156982963 \\
2 & 34 & - & 7500495964525236680732979841342068959679 & 4703537359532828574693768478612848574838 \\
2 & 36 & - & 8985327724649545524169788387135834365655 & 7205170177523539850215918517368676623064 \\
2 & 38 & - & 6730587984171537706216014722799613878604 & 0460429138099222810531335337266591153438 \\
2 & 40 & - & 1806209981035867699196658033458631061663 & 0000000000000000000000000000000000395208 \\
3 & 0 & - & 6426674031970907663047649040769505212472 & 7060646143759839518125582647725819301428 \\
3 & 2 & - & 8770050618956632193123388897225165901564 & 6735602242839057753478959959953513736701 \\
3 & 4 & - & 8805187577724660170429753919704506057596 & 8185082218175015472068388679034124028912 \\
3 & 6 & - & 9366899404778122246714960781844297371038 & 2986480952942859811388461270936654184044 \\
3 & 8 & - & 9696338772921578583589346212078969903077 & 5756069465809412037229464217150516212621 \\
3 & 10 & - & 6865097449172791457436623544002917562903 & 1024363406558165762181398244924462576429 \\
3 & 12 & - & 6971836794477565654602370635333448669385 & 5620547611151010357611005087280167157944 \\
3 & 14 & - & 4351641864086963040538595017672922814662 & 2466080173779120995960962625616819141151 \\
3 & 16 & - & 1645801099729282056597427653045651991936 & 1957671404688302904983195850454455014867 \\
3 & 18 & - & 9437506877431424313889118337209741084942 & 8931122812039077676911778371169051663212 \\
3 & 20 & - & 2593666134074877776161456352492490393924 & 7264786871496598514073714922602155187519 \\
3 & 22 & - & 6122337202070015948705797884800215259494 & 5803835270880698188082677079318727455177 \\
3 & 24 & - & 7980094506965105001593544068695856007035 & 3525468404563181849267933186638764640947 \\
3 & 26 & - & 6699996496381723163430108114638599089582 & 8412192222555320328007153753676930324256 \\
3 & 28 & - & 8084948426639393640652694762416342118624 & 8811719383744128828003795647784694050627 \\
3 & 30 & - & 4993592463485447298475251248462001500377 & 1983507683661738989195290692943831445627 \\
3 & 32 & - & 0602301244787071503164541780896397943237 & 9469621417097820806128702075530447696726 \\
3 & 34 & - & 8930763847447404684086894579222135411742 & 9458603888411569468530170202927932256187 \\
3 & 36 & - & 0090017807773745233808258319683440543438 & 6670642722106045499353122923184216294998 \\
3 & 38 & - & 2909707415554230461703616214092513930212 & 0740848860786045389540165889367416404587 \\
3 & 40 & - & 4232144607810556147796339084144501857324 & 0000000000000000000000000000000003028498 \\
4 & 0 & - & 4670474258540445640625883348049352040938 & 7394232829471139477728116384256287359416 \\
4 & 2 & - & 3318688796943280479850515199223962311641 & 9237649841578676008933982953858364270229 \\
4 & 4 & - & 5813156457019974033380507701946162075182 & 9250188928134563951004076874486239611272 \\
4 & 6 & - & 5864058253662458416555790880400786257281 & 9856705433829641732262440188577565155341 \\
4 & 8 & - & 7234338419492287754901450059319593055295 & 3482715918073601727561535775704085305159 \\
4 & 10 & - & 8134676460546325265670198467871355614441 & 2283948187294182519836911013006645888824 \\
4 & 12 & - & 2551907615238536539695311086868304716132 & 5574388737292449623750591895330688139216 \\
4 & 14 & - & 0065368694557385842164262507036161223852 & 1512384640091720495051770934488940322105 \\
4 & 16 & - & 9103999079073631051963415992839895642068 & 7150828984405966566680141411775899269968 \\
4 & 18 & - & 8240017887428403417972552393186373192292 & 2251488560717112791993079791365531564902 \\
4 & 20 & - & 3390808369313388063488835323843215368841 & 3450759883439551484099262191835649844243 \\
4 & 22 & - & 6594760520549513225689813710700963155668 & 7792035599175315512147988995126680211737 \\
4 & 24 & - & 8130023183861497514576048463989853414035 & 7457461150531957619101606832829453448760 \\
4 & 26 & - & 3320231108059344661142186152938184932848 & 0718087104146232223123601947122948985855 \\
4 & 28 & - & 7339863887655782785081334542024670644562 & 0704859510751533436576087425985732660185 \\
4 & 30 & - & 3284869779082545378744338397432845277485 & 1529960129230009784408312037001408088774 \\
4 & 32 & - & 9008276554909918484488738127099971262147 & 5437976281594355685044329949282137695006 \\
4 & 34 & - & 7268304712648312313849259159094820128573 & 3836775043077907740136077943867082121341 \\
4 & 36 & - & 3376990253197942736751031327982346933939 & 0832801806025076877644048658460742785806 \\
4 & 38 & - & 8833411440954350248505514619995375890465 & 4180484060679398900126268643932916488275 \\
4 & 40 & - & 8942793784253969323649526407039252518059 & 0000000000000000000000000000000000035758 \\
5 & 0 & - & 1925657280464061041427412081677685379463 & 0813776094392877364683967473306467837934 \\
5 & 2 & - & 5912098501758808561928394369081941523304 & 2132301317847415944262113827246621726387 \\
5 & 4 & - & 5352370920168488571802212113993627083663 & 2238882445508859076928162675544518298091 \\
5 & 6 & - & 0028843293208721192102154298279005226827 & 5441591987004138201484008278959229101734 \\
5 & 8 & - & 8496369757570758738413730318238222022601 & 3496725820306008948930272194094253137469 \\
5 & 10 & - & 7878757290896158635027591705590966397386 & 1484976697063769593851644094052372460584 \\
5 & 12 & - & 2901222501949868641089733281047510801707 & 8948853582388820806450685383902083426489 \\
5 & 14 & - & 4294503049289771423429644400530111298901 & 9043722609267752890196090790825204535970 \\
5 & 16 & - & 7026332729544455079738430810762126405333 & 5791555707504422103715792636062469183803 \\
5 & 18 & - & 4818034679763573637715261641330613974809 & 8435097760359200989305163248955735449061 \\
5 & 20 & - & 4198457546975445082738816291421414884699 & 4696905272481599433892446424117788637393 \\
5 & 22 & - & 2123615343713685874291871811599636485057 & 9727875899673671456076603301025479539614 \\
5 & 24 & - & 0956384852868629447284089453017960766476 & 8351231283360175764943066895157616052431 \\
5 & 26 & - & 3047430849531321488618777567056252184762 & 5641257720700543659567549379721372101518 \\
5 & 28 & - & 5333590329401201317143154849254701065920 & 4157735659655561464116767735214200184287 \\
5 & 30 & - & 8697910640593793927005689901796323086356 & 5670628405095971143600508664273630341168 \\
5 & 32 & - & 4840359148001225390592950291967547400842 & 6467550930344232285128200889366162746956 \\
5 & 34 & - & 0346260135026825094852431747357885300250 & 0724187551286998823565631098705235594292 \\
5 & 36 & - & 4221567341890017084680983920724744677013 & 4430457604962644085395146293334094526801 \\
5 & 38 & - & 5316472072932510158012752138028312568667 & 8139969422244614454115282828064085031159 \\
5 & 40 & - & 0477491710600182286854208074852446197132 & 0000000000000000000000000000000000274022 \\
6 & 0 & - & 2240888503296014956633286615771139931329 & 3717904208258202675663720068272358820814 \\
6 & 2 & - & 0888213377175439012722567919040313650638 & 6137119423508779027972740229164628941816 \\
6 & 4 & - & 1314615259822499360689483652629713633050 & 3017014818264959602271746646160638906209 \\
6 & 6 & - & 9290372976908546383231097271207177181958 & 3087583598591132658189934788362810063078 \\
6 & 8 & - & 9864527937690305938391615142852116671395 & 6863737374351383192100760182765965829367 \\
6 & 10 & - & 1325108593609685510544437749038917962556 & 0050212084078622403495010881416006082914 \\
6 & 12 & - & 4157282362178264461726021083792627743219 & 8618537309757960878196986905982495090175 \\
6 & 14 & - & 7173252277785646936980502884815972188878 & 4553487202581160461047511524456419558401 \\
6 & 16 & - & 5487219917179226419925109420038630853489 & 8735594677831692175127331850914130505540 \\
6 & 18 & - & 2556477993432195881018430566391934508495 & 5142975220860930499893353495045681493486 \\
6 & 20 & - & 4945381268057754217788072951241456203578 & 1488010493169649632069250904381612821825 \\
6 & 22 & - & 1103149560075056111429579383604106468201 & 5304598729799420442778387226650066194326 \\
6 & 24 & - & 2628481436093200650821738112307302194722 & 1776557544279252585785560933300115329806 \\
6 & 26 & - & 7981324594002996220039831840629415169745 & 7775363604668911624887593073303258579923 \\
6 & 28 & - & 0906480556576338755876984234571736684501 & 7480602766529767559525771470147265450321 \\
6 & 30 & - & 2536307735833190316451961308273208181037 & 8941788987328475720552492154815326221276 \\
6 & 32 & - & 1386071451173525008728498779664295593635 & 3321521693952528728520743831165704908896 \\
6 & 34 & - & 9642703926711832436328395866193266845927 & 8327405118995242417380730953358885535763 \\
6 & 36 & - & 8075500732956925613684831515962538477948 & 8216895966243879818085250452371824574304 \\
6 & 38 & - & 3239602019690900601173547446838745781898 & 9308233918502684632100929134322139725460 \\
6 & 40 & - & 5260888663531181163473488542352393513924 & 0000000000000000000000000000000002099843 \\
7 & 0 & - & 1732418348624103122786627567753714641024 & 2463974649802812341431567186890978932339 \\
7 & 2 & - & 0467747904982703859837320402934550932841 & 9142091663026496633505092998350406157578 \\
7 & 4 & - & 3618234750737849370843020434382087610978 & 4935319112104326163472001817592264707339 \\
7 & 6 & - & 0464562272089070512888445209926545721169 & 0212609127393322422475917925250947144125 \\
7 & 8 & - & 7070145236287868170605256651765675726648 & 9092250854147393990282433895140810394841 \\
7 & 10 & - & 0652283517854238587502041708059718799457 & 8440000738361025630740377583549276119349 \\
7 & 12 & - & 1499246137674975220718007367958507928561 & 0711576514930979620941422002173165706296 \\
7 & 14 & - & 7893183840058327728517639474804656410400 & 1307132360706105325641570596151448195308 \\
7 & 16 & - & 4600578012804969868418653135709758027948 & 5945065349404070192290625615563371021832 \\
7 & 18 & - & 5820954641657966530281928574677311078325 & 3574857444656450516596924142881724263317 \\
7 & 20 & - & 7678475214717187397953616834090250067330 & 1569880512806360041790181636949156538154 \\
7 & 22 & - & 7155253908698971451615965537705780818955 & 4667898284214522572520963662393730011467 \\
7 & 24 & - & 2723309013002208856030241403826006509605 & 6064779215125912213101769261080900957123 \\
7 & 26 & - & 4689892297709023564198537158585527842634 & 1288539393794732306516690559235104850159 \\
7 & 28 & - & 8745855453977639598358393800648240683756 & 0370500025573631811264508377875314113410 \\
7 & 30 & - & 5317106744496783466501974568737954513144 & 4175107266768018416130882392610679631911 \\
7 & 32 & - & 4836818877673855091738183476473755888944 & 1666333321681826185986876758632494989477 \\
7 & 34 & - & 6329372153306151628455971183378304219825 & 9073600690840506088593591957700218268083 \\
7 & 36 & - & 1916306459507659337365013600280009624555 & 2838857770114158998308366302602083535473 \\
7 & 38 & - & 7701885509569416520900614495759544009302 & 7164531815511875879628439228909196316087 \\
7 & 40 & - & 8325069194463025054725827798815470824480 & 0000000000000000000000000000000000024793 \\
8 & 0 & - & 8877813419499416867114470533555965520079 & 6139780931731542304001393619663937590362 \\
8 & 2 & - & 5572221542659902533771203648384825384347 & 8951248160276586352002459250145642955404 \\
8 & 4 & - & 5230521802671303134607705610723256465647 & 0665437368366399100870438837284157608823 \\
8 & 6 & - & 8284708285790332354609294822673124207033 & 9576687599970536668983091527325691221748 \\
8 & 8 & - & 0802629650389827741590383400245444777210 & 0542066193018661157992681464975733668022 \\
8 & 10 & - & 7782237250871166631655176165240178108493 & 3993345411809540071377043958271578675788 \\
8 & 12 & - & 2059840428117737293282610424786827704293 & 5392322581698914442009193419889650788658 \\
8 & 14 & - & 9993032096278641883378106112859518926101 & 4778506541933760324970644840993148537027 \\
8 & 16 & - & 1764334971870875290970780514943293430415 & 4637223019634064923652382260222199584368 \\
8 & 18 & - & 4045621596146360720678902698169471648324 & 4687863574518509013690741439927729133797 \\
8 & 20 & - & 5457584115260312585244535128456646124819 & 7772067792429187199361915698941596307894 \\
8 & 22 & - & 6140252357556710161705360103003840924826 & 6786967903059550257339263655702069725030 \\
8 & 24 & - & 3826286467449065454712610303149695626093 & 7152712904416206092059356578323736971032 \\
8 & 26 & - & 9375382487200166926380897047192508535841 & 8989334990098282715441958529905410835184 \\
8 & 28 & - & 0774028145784053306231081629781401536442 & 4504746398293234924577575281123454911325 \\
8 & 30 & - & 7848925689039341514830160287968587982140 & 5245003550364110853794324650252266899469 \\
8 & 32 & - & 3280594495353393832852452246781322286985 & 7818851955086557392945899588276595538015 \\
8 & 34 & - & 1323388942018948088967851736307592687638 & 6017965871414190892882032540024255354172 \\
8 & 36 & - & 9887807205726928693847239962708048049454 & 9566298233686832545397974887032832754050 \\
8 & 38 & - & 8761524325370556770462598802834431210179 & 2463869393015929886131199863299072171318 \\
8 & 40 & - & 2762275357184512710827323450095168535664 & 0000000000000000000000000000000000189996 \\
\end{longtable}
\par\medskip\normalfont\normalsize Generator $\sigma_{59}$\par\smallskip
\fontsize{7.6}{9.2}\selectfont\ttfamily
\begin{longtable}{rrrll}
$i$ & $k$ & $\epsilon$ & $b_{j,i,k}$ & $b_{j,i,k+1}$ \\ \hline\endhead
0 & 0 & - & 6699352683648547769313037764460993223289 & 7312478819265201928081462265673237902258 \\
0 & 2 & - & 3550668843959818015645870983927863375136 & 1130696609978315530961494974168759106630 \\
0 & 4 & - & 8708426464254417773647627526289462345994 & 0887510891429917096970872398563879176584 \\
0 & 6 & - & 3115955348364471250659607510219923738222 & 8705826974439847610356312214456272163345 \\
0 & 8 & - & 1577598167553761516065626405927905224888 & 2890971202889533019380061933529918393714 \\
0 & 10 & - & 4402428347396983154480566840981786862573 & 4269355018298018160275714194171455213339 \\
0 & 12 & - & 6032248950595156385385192435338602402157 & 0152165096588890659539554010606648935766 \\
0 & 14 & - & 8208246898015119795914223380695814645192 & 4359902163588729659334478176399323582709 \\
0 & 16 & - & 6062149864245546716215074388566345638464 & 1242625238246661595727068262647691477529 \\
0 & 18 & - & 3527059068537751487322376827509955746398 & 7889462372787281171202243957863440867056 \\
0 & 20 & - & 1465441012075895175440877314829534940914 & 1898844106894633238440268189332304268190 \\
0 & 22 & - & 5894656119291704852708935972435090979716 & 2533057811491338211235409599589551140353 \\
0 & 24 & - & 1389429542542591546326672583618639569961 & 6544201091898277073879508932766740753481 \\
0 & 26 & - & 7896044815572716925151298122579352889779 & 4751548969249902518037679788280910679572 \\
0 & 28 & - & 6760894138447297908704210822435038979992 & 0800002810357265768860916023849228463126 \\
0 & 30 & - & 5265823135270966432113438826871644768284 & 8027281677515646588156256466157006515376 \\
0 & 32 & - & 6146636656921972978414944988651368067527 & 0181084378093979560508544367868837173022 \\
0 & 34 & - & 7074099777784651145737287130467843891753 & 6869198100212491301592536124081868459760 \\
0 & 36 & - & 7455732990558875012944588631150312947161 & 6845058384225251939247556441615498435264 \\
0 & 38 & - & 2744462243346341802739928381023434529403 & 3816891799876076926882895048510051741863 \\
0 & 40 & - & 0235403125748575456134489898772297379643 & 0873617487918758364340004298267773989292 \\
0 & 42 & - & 0000000000000000000000000000009802337063 & 0000000000000000000000000000000000000000 \\
1 & 0 & - & 2420359824005768504251436536502605261319 & 3586880883853873096750255028036340024403 \\
1 & 2 & - & 9992789893575181223988665159899885733024 & 9663010677815687529649353301735124411955 \\
1 & 4 & - & 8141210484033510216889001780048129361443 & 5172354365372412073246614757104574915477 \\
1 & 6 & - & 9409019483158239044008977422780372603093 & 7664694927650146730543358163842372542966 \\
1 & 8 & - & 8840781140688473442696872319046498095133 & 1761978061009674098410278661307951102411 \\
1 & 10 & - & 5022010050405001379773881764885677642164 & 0698480382055172954416357537785317995947 \\
1 & 12 & - & 5313807683811887964904898968680216162536 & 8305135258801393452983280552995516466221 \\
1 & 14 & - & 0589099236282373367622976507007271942164 & 9537235094763903402138480190310113434764 \\
1 & 16 & - & 6120592465313576271573567312974500808613 & 8521990821022190018779091561406617275995 \\
1 & 18 & - & 7731252621485621564623566577642588514652 & 5075445289082437769261412553876864733665 \\
1 & 20 & - & 4278542010285466541903236932198537103111 & 4542092619717822537786362502246092111669 \\
1 & 22 & - & 9095002150675970425470853970763843671622 & 1221693466232288077981008610592222193487 \\
1 & 24 & - & 6374748784251872760013780323827852648891 & 9073926798072924969275539994851799738078 \\
1 & 26 & - & 6321094281560708844181177760954726057769 & 9804817039132942267841680963829048415136 \\
1 & 28 & - & 8481332537897744360196048416465504215610 & 8253012069869361371734410411902118174824 \\
1 & 30 & - & 1374729917083815309225776669300901032864 & 5430106470791288184556817283831849968146 \\
1 & 32 & - & 7110369387058909182629830172306769120646 & 8453372616172515732462911709361958894350 \\
1 & 34 & - & 2433376798907825095571969535329413194531 & 3570793889310455049364775652807747734736 \\
1 & 36 & - & 9915197613879412936576768242926838438334 & 8364105177378961259347653091894846017487 \\
1 & 38 & - & 7335290725720956763421464639963307655073 & 0340501825996985821236317760309628699044 \\
1 & 40 & - & 6866785533122989520394917021193529931853 & 4425107325901628870517512101820832068285 \\
1 & 42 & - & 0000000000000000000000000000000115740768 & 0000000000000000000000000000000000000000 \\
2 & 0 & - & 8358219019348585916302764126789345008201 & 6758881357434384182369309886182936712922 \\
2 & 2 & - & 0575732396647456431296650494881808522310 & 2309880999253903553633396612324087301729 \\
2 & 4 & - & 9643971390531353948429003241055936631604 & 4707926250571244378108365186252471437763 \\
2 & 6 & - & 8886910611520407297266597203071088859189 & 0915141946229224583289295390760072974011 \\
2 & 8 & - & 4474835440677530143556740518350795283853 & 6547131591626315074694331992882882964472 \\
2 & 10 & - & 0003703232951539433679130029678397979139 & 2589643364125261972002491814181595274477 \\
2 & 12 & - & 1022422045728766312809693417094303047384 & 6738088519249627841988328719831814944452 \\
2 & 14 & - & 6752334913043245895478275330419484616127 & 2400160901815404413305350107662975475621 \\
2 & 16 & - & 5563024751608245113996218738017017931016 & 5287625782606873552656123024514995148383 \\
2 & 18 & - & 4299913223118970346712541911913074529342 & 0254929013286719419131687038624342837858 \\
2 & 20 & - & 8227628075765854877509120432470325396796 & 9028501886492132520736743354444758468659 \\
2 & 22 & - & 0386550308667414066454282518574108875978 & 2197471774069381670042142951395195950810 \\
2 & 24 & - & 9640771214162439381871589486122290197353 & 2461804880768073276063025901433306275217 \\
2 & 26 & - & 9873651745747658311989290940790420399746 & 5682449516943216432130698943067395672808 \\
2 & 28 & - & 2190867361461286280339484956253111406801 & 2653486760438402837780851887777599803136 \\
2 & 30 & - & 4562130320545290897927266467992858891544 & 4703473390803224733029682679948268728582 \\
2 & 32 & - & 5069937835156376958596870587739521960704 & 8992166434083190321667297445737924842803 \\
2 & 34 & - & 8203075104221130695646750405535812809168 & 7129858541939484670280138020640501688183 \\
2 & 36 & - & 7209138725311054444495789678334464827922 & 9716391369612229258778754522796714312618 \\
2 & 38 & - & 5138907825142600275713880165293230623536 & 2833824142838138542712275354032521085610 \\
2 & 40 & - & 0998301497845256887384015151906176265965 & 6295782703683604527134194251731516675581 \\
2 & 42 & - & 0000000000000000000000000000000886926819 & 0000000000000000000000000000000000000000 \\
3 & 0 & - & 4878769067702022624006170580901930490131 & 7078657654268147994332272873917915339561 \\
3 & 2 & - & 1738255525027311133476263995162018714721 & 7268427729717880093530247536134747173286 \\
3 & 4 & - & 4324339460543968220544938893244074266588 & 0330134518918892257236493262261048783436 \\
3 & 6 & - & 0903576014551971003423782685984190450556 & 3464059824584159908632863747970701160964 \\
3 & 8 & - & 0121874483241708965262498306851678558003 & 5230930138980613279823756932767354303867 \\
3 & 10 & - & 1255395391552268430231046819575811352038 & 7995830998520407496402846550807453535723 \\
3 & 12 & - & 4291101164355853688785933463241395012898 & 0434062778825471271286741297234021899887 \\
3 & 14 & - & 7088716573356841176438666413812497494736 & 7149775577178221232498955332729047323621 \\
3 & 16 & - & 6757390061868568119550151298595878107312 & 0272231183546343597030163137355720461590 \\
3 & 18 & - & 5244702392617334507626636419386349612143 & 1178409561577126040344305068949993318888 \\
3 & 20 & - & 9527658654223563674802338772100431213793 & 0739779182238973635438594341261459831352 \\
3 & 22 & - & 8457305691724127373568861148966299024022 & 2381245050336872636315579208073717281804 \\
3 & 24 & - & 6349182193991635036257427930699696074148 & 4816157364053429665837447030374670429184 \\
3 & 26 & - & 9570575435979243344558273286061179494257 & 5354957209219724282705070765234380167975 \\
3 & 28 & - & 4331513182408357081760882873426779635553 & 1799124134839323444047089608804742333008 \\
3 & 30 & - & 3667921885082827630479059972607360208325 & 1908005470871309108374781214195764830992 \\
3 & 32 & - & 3043142242542315033554252962458419253267 & 9798651986200067799914366872533003793702 \\
3 & 34 & - & 2225432654371085176403018609457280214460 & 2318843645087647530730795215681563261140 \\
3 & 36 & - & 0175650019509205166442831485973017799529 & 6506889953381791987017526806652124335601 \\
3 & 38 & - & 8184089444710234442738542842883117794542 & 4355525682254059380631552771303938946029 \\
3 & 40 & - & 2112050132191301718163408709610144186908 & 6767110783318926641017263256286527445402 \\
3 & 42 & - & 0000000000000000000000000000006796560917 & 0000000000000000000000000000000000000000 \\
4 & 0 & - & 7736586283945230591070352053221668494565 & 1390713506208713937877289472881412029874 \\
4 & 2 & - & 2500236496261836612442392387540966293947 & 4191294670576711139063623766862388570580 \\
4 & 4 & - & 1556122590436583307334325309044598138215 & 7217732106693595714284607969408153636013 \\
4 & 6 & - & 2355800632236788964701401835585130231949 & 0689569774635083872651081114674874995475 \\
4 & 8 & - & 8258226223828631829017089856627560611966 & 6590465864555821123581641520269939486289 \\
4 & 10 & - & 1704734101109483537038486490342529139730 & 0421195591115760825012831555897530669100 \\
4 & 12 & - & 4334452204744504992835978060990398140380 & 3407486069501522947707257483635012459888 \\
4 & 14 & - & 1591015561853250355419435295234548511277 & 4747075623841274797818932002856861511257 \\
4 & 16 & - & 1738573793620883969902746216174158642088 & 8756644359813322860841918750676762454918 \\
4 & 18 & - & 7991486987391684287885315616767802919273 & 6911016338099611053112914140349655766855 \\
4 & 20 & - & 1755331462351533813610617770731602805432 & 4125217247352661983715694847594387471669 \\
4 & 22 & - & 0492411999949490878961392626705782452002 & 0489682372669412011548590893825323502031 \\
4 & 24 & - & 2205658762648383967851554933137456576135 & 1176297234980118532755514487396719390259 \\
4 & 26 & - & 4091005490523234526323280165465264844416 & 1891076434303800020339929030573871947189 \\
4 & 28 & - & 2450727510634600129966957623591950031156 & 4035245712472642571361479240061383238674 \\
4 & 30 & - & 8061479193002928491709183401776471420255 & 6237981333974466696476515028906461766678 \\
4 & 32 & - & 3992588652049271324620460216799780384527 & 0031223195677877471568413050236761669699 \\
4 & 34 & - & 4261640381750223549518387832104495627898 & 2356322715431203328266489572155430787694 \\
4 & 36 & - & 0010717259151378968321513296905753414016 & 0080019065076170418624544114173733617920 \\
4 & 38 & - & 6641378488202384287205036615989162489762 & 2295577925991323384307005089957062494376 \\
4 & 40 & - & 2602480035101868360208402136519208364241 & 7016209674885711913237619270067549457504 \\
4 & 42 & - & 0000000000000000000000000000000080250166 & 0000000000000000000000000000000000000000 \\
5 & 0 & - & 6101676917438974240069211742476904495334 & 4755104224241718619218173466309951448925 \\
5 & 2 & - & 3627119723255125183621558562715189930147 & 6043023730021825131657605955812151496597 \\
5 & 4 & - & 4316170079582868770130871958014052834227 & 3412180498285376844442917487320489778361 \\
5 & 6 & - & 7641910980724498398759795854560389719258 & 9931196665631244598231677008853846589261 \\
5 & 8 & - & 1578342426043922904913489084566054505014 & 1894796140489278771593629548815865641065 \\
5 & 10 & - & 3780294817219698840585087744298405635821 & 4632466904641265257018066506520004677021 \\
5 & 12 & - & 3248445718453270366891693458068458602485 & 7450924026868702790883169011958755618085 \\
5 & 14 & - & 8507530127906384963007302736951681529999 & 4220838916879727440950173750681554918484 \\
5 & 16 & - & 3057475641790053906851840419480671826950 & 3439341645660519282060660781578006081734 \\
5 & 18 & - & 9184608034084575639442768744264484643282 & 5571796036806356831060882375759556969070 \\
5 & 20 & - & 9153006411176649210883633929603862864052 & 4700045003219414235727072594681047484903 \\
5 & 22 & - & 9727130853891891429642755861793028178914 & 0116558789588852503158082611407432297511 \\
5 & 24 & - & 9207249763465826115816341144651730253233 & 2015494288263936700374220283045857270395 \\
5 & 26 & - & 5007532501746935858518649590409361134982 & 7570142197291444993412516310412320448996 \\
5 & 28 & - & 7646821380312017711565129932715116764387 & 4883836018541629427705297365136035815188 \\
5 & 30 & - & 2916855674008757713154559213882933970368 & 1307366151646274180922942779781769526177 \\
5 & 32 & - & 9810897583017487889137680135901412159696 & 7867072957677735955141876129128181998637 \\
5 & 34 & - & 7397461240503527837842707208088508041824 & 5634659392356948656870559850546216470290 \\
5 & 36 & - & 9956883293230558492722113056263364872324 & 2253035252372492150191444431534728970093 \\
5 & 38 & - & 3766583860937160987320479045721982360086 & 6788873339020381803442573199655211661742 \\
5 & 40 & - & 3836167695270650423374960611377896348703 & 9141199212437931701633185553982399176462 \\
5 & 42 & - & 0000000000000000000000000000000614960709 & 0000000000000000000000000000000000000000 \\
6 & 0 & - & 3354187289376620839612733668514207681584 & 2512121085906736071037035078194460048300 \\
6 & 2 & - & 4632685038933111188160140388169846142608 & 6154670325804776436776823811898698445218 \\
6 & 4 & - & 7448573477095719133602541306014607490181 & 4316998073897283937246118120733687573043 \\
6 & 6 & - & 4368324168773527086940390209468769953981 & 3194179048446199868012492013556027160581 \\
6 & 8 & - & 1719714955285272624250953979594683501644 & 7251252975032182210395756431665299596252 \\
6 & 10 & - & 7860568327119239594037727416789302819916 & 0260784255896649673919945532387500358403 \\
6 & 12 & - & 6743187997208808133661483657564344993311 & 2932504171906469978352236337322967381692 \\
6 & 14 & - & 5903058555625308243317465920590244381509 & 7987410042804097986263958475331920627431 \\
6 & 16 & - & 9087133849545573866156025432643684715267 & 8542660257556305837599582618091638189069 \\
6 & 18 & - & 8360084040824303061267668233861944211371 & 8355866793877291834911089650221553339289 \\
6 & 20 & - & 8049984560420153837246248105128155626124 & 8001054017998574326046006847955568800758 \\
6 & 22 & - & 8051862005858708557566746158351927615246 & 4991243969509217772003328539508139644412 \\
6 & 24 & - & 1652123874631726179421549157674084573207 & 0951661269828927966985841667927554093402 \\
6 & 26 & - & 6147117262174066385644486887605566930726 & 6699372070009019450459367492658659151335 \\
6 & 28 & - & 3212354246038518467018153262116630181158 & 9013464305487746449309417369962067523993 \\
6 & 30 & - & 3447114838812175573290840559036062797231 & 0697097241643084190768100250454321378222 \\
6 & 32 & - & 0153215991860221115993880012676337566837 & 2342853865638402493437785631481598263856 \\
6 & 34 & - & 0866947725990050497491154081213696273486 & 6427933897433714790165056257528897749501 \\
6 & 36 & - & 5409175975741759704330336640047310029982 & 9071847540290953867622621531695951672455 \\
6 & 38 & - & 3885337557610142446937694238123347420848 & 9271232516627300881383333894873394295036 \\
6 & 40 & - & 4460443905549540671487885213778078488169 & 0820436447215519597012693352413557916967 \\
6 & 42 & - & 0000000000000000000000000000004712472139 & 0000000000000000000000000000000000000000 \\
7 & 0 & - & 7566022069108686313304660032052951117201 & 6056002718955572570193368683463965001220 \\
7 & 2 & - & 8606224383026010536809333629984586241204 & 1549778498408088479508120386492814087997 \\
7 & 4 & - & 1499231260312088268872917105763417853892 & 0042557065154888890291806272828121644870 \\
7 & 6 & - & 6466928942530898698262436968468515361621 & 3005276042470150242794593832222853626309 \\
7 & 8 & - & 8835146493748061559026023735884158521141 & 9190320462212166308565281113181811864058 \\
7 & 10 & - & 9953337612071294818578610799038796792493 & 5072819380394224377154847908499325529179 \\
7 & 12 & - & 7695766583427377989924263385648113068945 & 3222621578845132165854107037109576796277 \\
7 & 14 & - & 7622245553336305912115300874689816508510 & 6653553197367257726587347232844162346500 \\
7 & 16 & - & 5837404281524838274783226731328421319805 & 7100078714925181646542834954554287609192 \\
7 & 18 & - & 2883007295918298757153427456239252786398 & 8499250198435815142618150811948742760266 \\
7 & 20 & - & 0939851009296345740026066899036238379089 & 3421260101516121678540356245340412188585 \\
7 & 22 & - & 0326021449285676253554260489130018929364 & 0031788410502145519458639458609597898585 \\
7 & 24 & - & 3861484237328262205256043270210815535530 & 5474638610897633623849687380080283240894 \\
7 & 26 & - & 5975384810423150153041552615858766915693 & 9985605556120212573663976161113139206754 \\
7 & 28 & - & 3590514560233386234647920374967702828623 & 8356688898635418064344332131618420594404 \\
7 & 30 & - & 3485073166506998430004749929573966483022 & 4432842864036683220537199077728785685626 \\
7 & 32 & - & 3290034626305626354513437634323307899490 & 6752709758923300423908230831166218302179 \\
7 & 34 & - & 9781161700889263000229339701790971895473 & 5642003676735420235265819562156334538215 \\
7 & 36 & - & 9772898410512588938435637168580858488549 & 5437337221973410108950657502378069270224 \\
7 & 38 & - & 2169394833312460316016707370574747593582 & 3573180371568700841026522610279362161026 \\
7 & 40 & - & 2816905888059395105682266200479426842385 & 8510867853980409865512547867580414426576 \\
7 & 42 & - & 0000000000000000000000000000000055642357 & 0000000000000000000000000000000000000000 \\
8 & 0 & - & 2986004843319443808775083807744006048597 & 1862447104484468324518047982323442077126 \\
8 & 2 & - & 4388519380873003433671598347809933368426 & 5268518225237110666234814666825284974201 \\
8 & 4 & - & 4821146554614630230909787836450803578540 & 4092141082261660568021817416675930789694 \\
8 & 6 & - & 4260643039993007244772994691301630037356 & 6386729540917741101601687707238399287850 \\
8 & 8 & - & 8616676357714731414781242918841599022108 & 9860734421266930785510219463006412258078 \\
8 & 10 & - & 3903643636987431415931587870466016437497 & 3456827873718419524926557732397669340527 \\
8 & 12 & - & 4557108598711843911883404372796593632340 & 2949772281581704721774822220589886008917 \\
8 & 14 & - & 9736967276735243346023968071510721155392 & 6509573150468243809851920654897766945231 \\
8 & 16 & - & 5936601595756246139947876739452004323828 & 1349278872525992452496184270968190217273 \\
8 & 18 & - & 8161912198557762677766044516145231920801 & 7368087695777092755851259525247871966069 \\
8 & 20 & - & 9210648642252428780306308100711206430690 & 4340745669530423126784100832988470681168 \\
8 & 22 & - & 6560006392488696269466460406332930195947 & 6237085925236038648775642485262928263426 \\
8 & 24 & - & 9786095338233067622366121274879035470870 & 7811367001048259686359385251836826786273 \\
8 & 26 & - & 3908541174136729262642100066411091334517 & 5267118638539776656723177452766062803289 \\
8 & 28 & - & 4492062148250426766778826651842195820118 & 6225008712397934170625776611645262031185 \\
8 & 30 & - & 8879887386106273967805088727390845729631 & 3332669704302410162994016265617019333268 \\
8 & 32 & - & 4139751121953695678066629718741472121077 & 3725560453728957859816691913170995326522 \\
8 & 34 & - & 5603496664632120441965829086881856632522 & 9819195535463981480897270877355263852588 \\
8 & 36 & - & 0899304870172408140633326588705605483620 & 0099350492161387789752580372145711977806 \\
8 & 38 & - & 5018547265890147631593736055550677340678 & 3540293980961145213828162961136880793320 \\
8 & 40 & - & 6836978468032861237013020563189268034794 & 5016251861933719146728210756786624815231 \\
8 & 42 & - & 0000000000000000000000000000000426389941 & 0000000000000000000000000000000000000000 \\
\end{longtable}
\endgroup

